
En
posant \(y=1-\exp(-x)\), on se ramène à calculer le résidu de :
\[
 \frac{(1-y)^{-k-1}\cdot dy}{y^{r+1}}.
\]
c'est-à-dire le coefficient de \(y^r\) dans \((1-y)^{-k-1}\). Pour \(-r\leq k<0\), il est nul, et pour \(k=0\), c'est \(1\).

Ainsi s'achève la démonstration du théorème de Riemann--Roch.

\subsection*{BIBLIOGRAPHIE}
\begin{enumerate}[label={[\arabic*]}]
\item A. Borel et J.P. Serre, Le théorème de Riemann--Roch, Bull. Soc. Math. de France, t. 86, 1958.
\item A. Grothendieck, Techniques de construction et théorèmes d'existence en Géométrie Algébrique, Sém. Bourbaki, 1960-61.
\end{enumerate}




